% !PRINTED 178
\weberpage{178}
Eine große Klasse von Teilern der Gruppe $\mathfrak{L}$ mit endlichem
Index bilden die sogenannten \gesperrt{Kongruenzgruppen}, die, wenn
$m$ irgend eine ganze Zahl ist, durch die vier Kongruenzen
\[
\begin{pmatrix} \alpha, & \beta \\ \gamma, & \delta \end{pmatrix} \equiv \begin{pmatrix} 1, & 0 \\ 0, & 1 \end{pmatrix} (\mathrm{mod}\ m)
\]
charakterisiert sind. Die zu diesen Gruppen gehörigen Modulfunktionen
 heißen \gesperrt{Kongruenzmoduln} $m$ter Stufe. Für $m = 2$
ist dies die Gruppe $\mathfrak{A}$ (\S\ 36). Wir werden von den folgenden
Sätzen Gebrauch machen:

1. \gesperrt{Der Modul $\varkappa^2$ ist nach \S\ 52 eine zu dieser Gruppe gehörige}
 Modulfunktion; da nach \S\ 36 alle diese Transformationen
aus den beiden
\[
\begin{pmatrix} 1, & 0 \\ 2, & 1 \end{pmatrix}, \qquad \begin{pmatrix} 1, & 2 \\ 0, & 1 \end{pmatrix}
\]
zusammensetzbar sind, so können wir den Satz aussprechen:

\gesperrt{Jede Modulfunktion, die durch die beiden Vertauschungen
}
\[
(\omega,\ \omega+2), \qquad \left(\omega,\ \frac{\omega}{1+2\omega}\right)
\]
\gesperrt{ungeändert bleibt, ist eine rationale Funktion von $\varkappa^2$.}

2. Die Funktion $\varkappa^2\varkappa'^2$ gehört zu der aus der ersten und
zweiten Klasse des \S\ 36 gemeinschaftlich gebildeten Gruppe $\mathfrak{A}'$.
Diese Gruppe läßt sich durch Wiederholung der beiden Transformationen
\[
\begin{pmatrix} 1, & 0 \\ 2, & 1 \end{pmatrix}, \qquad \begin{pmatrix} 0, & 1 \\ -1, & 0 \end{pmatrix}
\]
ableiten, und wir können daher den Satz aussprechen:

\gesperrt{Jede Modulfunktion, die durch die beiden Vertauschungen
}
\[
(\omega,\ \omega+2), \qquad \left(\omega,\ -\frac{1}{\omega}\right)
\]
\gesperrt{ungeändert bleibt, ist eine rationale Funktion von $\varkappa^2\varkappa'^2$.}
Hieraus folgt noch:

3. \gesperrt{Jede Modulfunktion, die durch $(\omega,\omega+2)$ ungeändert
bleibt und durch $\left(\omega, -\dfrac{1}{\omega}\right)$ ihr Zeichen ändert, ist das
Produkt von $(\varkappa'^2 - \varkappa^2)$ mit einer rationalen Funktion
von $\varkappa^2\varkappa'^2$.}
